\chapter{Power, SDRAM electrical, and configuration}

\section{Power architecture -- status: OPEN}
Rail \emph{voltages} are real, DATASHEET VALUEs from the actual ECP5
and \code{AS4C4M16SA-6TIN} datasheets; regulator selection, current
budget, and decoupling are \textbf{not} finalized in this revision.

\begin{tabularx}{\textwidth}{L{2.6cm}L{2.2cm}Y}
\toprule
\rowh \thd{Rail} & \thd{Nominal} & \thd{Status} \\
\midrule
VCC (core) & 1.1\,V & DATASHEET VALUE; regulator TBD \\
\rowa VCCAUX & 2.5\,V & DATASHEET VALUE; regulator TBD \\
VCCIO 6/7 (SDRAM) & 3.3\,V (assumed) & matches SDRAM's own LVCMOS33 need; not independently re-verified \\
\rowa SDRAM VDD/VDDQ & 3.3\,V & DATASHEET VALUE (\code{AS4C4M16SA-6TIN}) \\
\bottomrule
\end{tabularx}

No real power-estimation tool was run against the actual synthesized
netlist this revision; a real, computed current budget therefore
remains an OPEN item, not an invented number.

\section{SDRAM electrical / timing sign-off -- status: PARTIAL}
\code{AS4C4M16SA-6TIN} is a standard JEDEC SDR SDRAM; the controller
(\code{sdram\_controller.v}) implements a real power-up wait, mode
register set, and periodic AUTO REFRESH, verified functionally in
simulation (init/refresh/read/write/burst/masked-write, 40 real
refresh events per D-Stress run, zero corruption). A byte-for-byte
cross-check of the controller's own timing constants against the
device's real datasheet parameters (CAS latency, \sig{tRCD},
\sig{tRP}, refresh interval, setup/hold) was \textbf{not} independently
re-performed this revision -- simulation-level correctness is
established; a dedicated datasheet-parameter audit remains an OPEN
item before board layout.

\section{Configuration -- status: OPEN}
Standard ECP5 JTAG (\sig{TDI}/\sig{TDO}/\sig{TCK}/\sig{TMS}) and
configuration (\sig{PROGRAMN}/\sig{INITN}/\sig{DONE}/\sig{CCLK}) pins
are identified and standard (real balls listed in
\code{docs/pinouts.md}); JTAG remains available regardless of boot
mode. No configuration-flash part number has been selected, and no
SPI-flash-boot-vs-JTAG-only decision has been made this revision.
